\documentclass[11pt]{article}
\usepackage[margin=18mm,headheight=22pt]{geometry}
\usepackage{amsmath}
\usepackage{fontspec}
\usepackage{unicode-math}
\setmainfont{NewCM10}[Extension=.otf,UprightFont=*-Book,BoldFont=*-Bold,ItalicFont=*-BookItalic,BoldItalicFont=*-BoldItalic]
\setsansfont{NewCMSans10}[Extension=.otf,UprightFont=*-Book,BoldFont=*-Bold,ItalicFont=*-BookOblique,BoldItalicFont=*-BoldOblique]
\setmathfont{NewCMMath-Book.otf}
\usepackage{xcolor}
\definecolor{FAUblue}{HTML}{04316A}
\definecolor{Paperblue}{HTML}{EEF2F7}
\definecolor{EEIgold}{HTML}{CC9C42}
\usepackage{etoolbox,enumitem,booktabs,tabularx,array,needspace,microtype}
\usepackage[most]{tcolorbox}
\usepackage{titlesec,fancyhdr}
\usepackage[hidelinks]{hyperref}
\usepackage{draftwatermark}
\SetWatermarkText{Made by Thanate Amornjiraporn (Focus)}
\SetWatermarkColor{black!7}\SetWatermarkFontSize{30pt}\SetWatermarkAngle{45}
\setlength{\parindent}{0pt}
\setlength{\parskip}{5pt}
\setlist{itemsep=2pt,topsep=3pt,leftmargin=1.6em,beginpenalty=10000,endpenalty=10000}
\raggedbottom
\clubpenalty=10000 \widowpenalty=10000 \displaywidowpenalty=10000 \interlinepenalty=10000
\setcounter{secnumdepth}{0}
\setcounter{tocdepth}{1}
\pagestyle{fancy}\fancyhf{}
\fancyhead[L]{\small\sffamily\color{FAUblue} From Fundamentals to AI Apocalypse · APDR · Chapter 9}
\fancyhead[R]{\small\sffamily Testing and Debugging}
\fancyfoot[C]{\small\thepage}
\titleformat{\section}{\Large\bfseries\sffamily\color{FAUblue}}{}{0pt}{}
\titleformat{\subsection}{\large\bfseries\sffamily\color{FAUblue}}{}{0pt}{}
\titlespacing*{\section}{0pt}{16pt}{6pt}
\preto\section{\Needspace{8\baselineskip}}
\preto\subsection{\Needspace{6\baselineskip}}
\newtcolorbox{equationcard}[1]{colback=Paperblue,colframe=Paperblue,boxrule=0pt,arc=0pt,left=10pt,right=10pt,top=6pt,bottom=6pt,title={#1},fonttitle=\small\bfseries\sffamily,coltitle=FAUblue,before skip=8pt,after skip=8pt}
\newtcolorbox{workedbox}[1]{colback=white,colframe=EEIgold,boxrule=0pt,leftrule=3pt,arc=0pt,left=10pt,top=4pt,bottom=4pt,title={#1},fonttitle=\small\bfseries\sffamily,coltitle=FAUblue,attach title to upper={\par},before skip=8pt,after skip=8pt}
\newtcolorbox{recallbox}[1]{colback=white,colframe=FAUblue,boxrule=0pt,leftrule=3pt,arc=0pt,left=10pt,top=4pt,bottom=4pt,title={#1},fonttitle=\small\bfseries\sffamily,coltitle=FAUblue,attach title to upper={\par},before skip=8pt,after skip=8pt}
\usepackage{environ}
\newsavebox\measurebox
% Typeset the section once invisibly to measure it; start a new page if it will not fit in what is left.
\makeatletter
\newif\ifkeep@measuring
\NewEnviron{keepsection}{%
  \ifkeep@measuring \BODY \else
    \edef\keep@sec{\the\value{section}}\edef\keep@sub{\the\value{subsection}}%
    \setbox\measurebox\vbox{\hsize\linewidth\keep@measuringtrue\def\Needspace##1{}\BODY}%
    \setcounter{section}{\keep@sec}\setcounter{subsection}{\keep@sub}%
    \dimen2=\ht\measurebox \advance\dimen2\dp\measurebox
    \dimen@=\textheight
    \ifdim\pagegoal<\maxdimen
      \dimen@=\pagegoal \advance\dimen@-\pagetotal
      \ifdim\dimen2>\dimen@ \clearpage \dimen@=\textheight \fi
    \fi
    % if it fits, nothing inside may force an early break
    \ifdim\dimen2>\dimen@ \BODY \else \begingroup\def\Needspace##1{}\BODY\endgroup \fi
  \fi}
\makeatother
\newtcolorbox{notebox}{colback=Paperblue,colframe=Paperblue,boxrule=0pt,arc=0pt,left=8pt,right=8pt,top=4pt,bottom=4pt,before skip=6pt,after skip=6pt}
\newtcolorbox{plainbox}{colback=white,colframe=FAUblue,boxrule=0pt,leftrule=3pt,arc=0pt,left=10pt,right=8pt,top=4pt,bottom=4pt,before skip=6pt,after skip=8pt}
\usetikzlibrary{arrows.meta}
\title{\sffamily\bfseries\color{FAUblue}Testing and Debugging}
\author{Algorithms, Programming, and Data Representation \ensuremath{\cdot} Module 23131 \ensuremath{\cdot} Chapter 9\\[2pt]\small From Fundamentals to AI Apocalypse · personal study guide}
\date{}
\begin{document}
\maketitle
\begingroup\footnotesize\setlength{\parskip}{0pt}\tableofcontents\endgroup
\clearpage

\Needspace{9\baselineskip}
{\large Specify what a program promises, design tests that challenge that promise, and investigate failures with evidence. Exceptions make runtime errors explicit; assertions check internal assumptions.}


\begin{keepsection}
\Needspace{13\baselineskip}
\section{Start from memory}

\Needspace{10\baselineskip}
Answer these before reading.

\Needspace{8\baselineskip}
\begin{enumerate}
\item How do testing and debugging differ?
\item Why test boundary values?
\item Does branch coverage prove correctness?
\item What does int("cat") raise?
\item Should assert validate user input?
\end{enumerate}

\nopagebreak[4]{\small\itshape Answers: see “Start from memory” under Answers at the end.}
\end{keepsection}

\begin{keepsection}
\Needspace{16\baselineskip}
\section{How to use this guide}

\Needspace{13\baselineskip}
For every test, write its input, expected result, and reason. Keep the smallest failing case when debugging.

\Needspace{10\baselineskip}
\begin{enumerate}
\item \textbf{Day 0:} §9.1–9.3; practice 1, 4, 7, 10, 13.
\item \textbf{Day 1:} §9.4–9.5; practice 2, 5, 8, 11, 14.
\item \textbf{Day 3:} §9.6–9.8; practice 3, 6, 9, 12, 15, 18.
\item \textbf{About day 7:} Rebuild the derivations; solve 16–20.
\item \textbf{About day 14:} Work both exam problems in about 20 minutes each.
\item \textbf{Before the exam:} Tick only skills you can demonstrate without notes.
\end{enumerate}

\begin{notebox}\textbf{When stuck:} reproduce the failure, observe values, test one hypothesis at a time.\end{notebox}
\end{keepsection}

\begin{keepsection}
\Needspace{15\baselineskip}
\section{How the chapter fits together}

\Needspace{12\baselineskip}
Specifications define expected behavior. Black-box tests choose representative inputs from the contract; glass-box tests inspect code paths. Debugging explains mismatches. Exceptions signal failures; assertions check internal invariants.

\Needspace{6\baselineskip}
\textbf{Specify}

\Needspace{4\baselineskip}
\begin{itemize}
\item \Needspace{4\baselineskip}
\textbf{Contract} \hfill{\small\color{FAUblue}§9.1}

\Needspace{4\baselineskip}
\texttt{input \ensuremath{\to} promised output}

Rests on: assumptions and guarantees.
\item \Needspace{4\baselineskip}
\textbf{Partitions} \hfill{\small\color{FAUblue}§9.3}

\Needspace{4\baselineskip}
\texttt{classes + boundaries}

Rests on: expected behavior.
\end{itemize}

\Needspace{6\baselineskip}
\textbf{Inspect}

\Needspace{4\baselineskip}
\begin{itemize}
\item \Needspace{4\baselineskip}
\textbf{Paths} \hfill{\small\color{FAUblue}§9.4}

\Needspace{4\baselineskip}
\texttt{each decision outcome}

Rests on: visible control flow.
\item \Needspace{4\baselineskip}
\textbf{Debug} \hfill{\small\color{FAUblue}§9.5}

\Needspace{4\baselineskip}
\texttt{observe \ensuremath{\to} hypothesize \ensuremath{\to} test}

Rests on: repeatable evidence.
\end{itemize}

\Needspace{6\baselineskip}
\textbf{Respond}

\Needspace{4\baselineskip}
\begin{itemize}
\item \Needspace{4\baselineskip}
\textbf{Exceptions} \hfill{\small\color{FAUblue}§9.6}

\Needspace{4\baselineskip}
\texttt{raise \ensuremath{\to} handler}

Rests on: a deliberate recovery policy.
\item \Needspace{4\baselineskip}
\textbf{Assertions} \hfill{\small\color{FAUblue}§9.8}

\Needspace{4\baselineskip}
\texttt{assert condition}

Rests on: an internal assumption.
\end{itemize}
\end{keepsection}

\begin{keepsection}
\Needspace{9\baselineskip}
\section{9.1 Defensive programming}

\Needspace{6\baselineskip}
\begin{plainbox}\textbf{In plain words:} make assumptions visible and each program part small enough to test.\end{plainbox}

\Needspace{10\baselineskip}
A function contract states acceptable inputs and promised output. Modularity makes parts testable in isolation. Remove syntax and static errors first, then compare chosen input/output pairs with the specification.

\begin{workedbox}{Worked example \ensuremath{\cdot} absolute value contract}
\Needspace{5\baselineskip}
\begin{enumerate}
\item Assume x is an integer.
\item Return x if x \ensuremath{\geq} 0 and \ensuremath{-}x otherwise.
\item Expected cases: 4\ensuremath{\to}4, \ensuremath{-}4\ensuremath{\to}4, 0\ensuremath{\to}0.
\end{enumerate}
\end{workedbox}

\Needspace{4\baselineskip}
\Needspace{4\baselineskip}
\begin{notebox}\textbf{Trap:} A function running once does not prove it meets its contract.\end{notebox}

\Needspace{8\baselineskip}
\begin{notebox}\textbf{Near miss:} Testing only x=4 misses a faulty zero boundary; add x=0.\end{notebox}

\begin{recallbox}{Recall \ensuremath{\cdot} §9.1}
\Needspace{4\baselineskip}
\begin{enumerate}
\item What does a contract state?
\item Why modularize?
\end{enumerate}

\nopagebreak[4]{\small\itshape Answers: see “9.1 Defensive programming” under Answers at the end.}
\end{recallbox}
\end{keepsection}

\begin{keepsection}
\Needspace{10\baselineskip}
\section{9.2 Test levels}

\Needspace{7\baselineskip}
\begin{plainbox}\textbf{In plain words:} test one unit, keep tests for old bugs, and check components together.\end{plainbox}

\Needspace{5\baselineskip}
\begin{itemize}
\item \textbf{Unit:} one function.
\item \textbf{Regression:} prevent a fixed bug from returning.
\item \textbf{Integration:} check the combined program.
\end{itemize}

\Needspace{5\baselineskip}
For each input, record an expected result before execution. Integration-only testing makes failures harder to localize.

\Needspace{4\baselineskip}
\Needspace{4\baselineskip}
\begin{notebox}\textbf{Trap:} A passing unit test does not prove the integrated report is correct.\end{notebox}

\Needspace{8\baselineskip}
\begin{notebox}\textbf{Near miss:} Two functions pass alone but use mismatched record shapes; integration exposes it.\end{notebox}

\begin{recallbox}{Recall \ensuremath{\cdot} §9.2}
\Needspace{4\baselineskip}
\begin{enumerate}
\item What does a regression test protect?
\item What does integration testing check?
\end{enumerate}

\nopagebreak[4]{\small\itshape Answers: see “9.2 Test levels” under Answers at the end.}
\end{recallbox}
\end{keepsection}

\begin{keepsection}
\Needspace{8\baselineskip}
\section{9.3 Black-box testing}

\Needspace{5\baselineskip}
\begin{plainbox}\textbf{In plain words:} choose cases from the specification without looking at implementation.\end{plainbox}

\Needspace{10\baselineskip}
Partition inputs into natural classes and test boundaries: empty/singleton, small/large, and either side of a decision. A square-root function might promise x\ensuremath{\geq}0, ε\ensuremath{>}0 and x\ensuremath{-}ε\ensuremath{\leq}r²\ensuremath{\leq}x+ε.

\begin{workedbox}{Worked example \ensuremath{\cdot} square-root cases}
\Needspace{8\baselineskip}
\begin{enumerate}
\item (0, 0.0001) checks zero.
\item (25, 0.0001) checks a perfect square.
\item (0.05, 0.0001) checks values below one.
\item (2, 0.0001) checks an irrational root against tolerance.
\item Very small and large values probe extremes.
\end{enumerate}
\end{workedbox}

\Needspace{4\baselineskip}
\Needspace{4\baselineskip}
\begin{notebox}\textbf{Trap:} Many typical values can still miss a boundary.\end{notebox}

\Needspace{8\baselineskip}
\begin{notebox}\textbf{Near miss:} Testing \ensuremath{-}2 and 2 covers branches in the faulty abs example but misses \ensuremath{-}1.\end{notebox}

\begin{recallbox}{Recall \ensuremath{\cdot} §9.3}
\Needspace{4\baselineskip}
\begin{enumerate}
\item Why test boundaries?
\item Why can another person design tests?
\end{enumerate}

\nopagebreak[4]{\small\itshape Answers: see “9.3 Black-box testing” under Answers at the end.}
\end{recallbox}
\end{keepsection}

\begin{keepsection}
\Needspace{11\baselineskip}
\section{9.4 Glass-box testing}

\Needspace{8\baselineskip}
\begin{plainbox}\textbf{In plain words:} inspect decisions in code and exercise their outcomes.\end{plainbox}

\begin{equationcard}{Coverage and correctness}
\[ \begin{gathered} \texttt{executed paths \ensuremath{\neq} correct behavior} \end{gathered} \]
{\small Pair path checks with boundary cases from the specification.}
\end{equationcard}

\Needspace{10\baselineskip}
Path-complete testing attempts every path, but loops have many lengths and coverage does not prove correctness. For faulty code \texttt{if x \ensuremath{<} -1: return -x}, tests 2 and \ensuremath{-}2 cover both branches, yet \ensuremath{-}1 returns \ensuremath{-}1, contrary to the absolute-value contract.

\begin{workedbox}{Worked example \ensuremath{\cdot} coverage blind spot}
\Needspace{7\baselineskip}
\begin{enumerate}
\item 2 takes the false branch and returns 2.
\item \ensuremath{-}2 takes the true branch and returns 2.
\item \ensuremath{-}1 takes the false branch and wrongly returns \ensuremath{-}1.
\item The boundary test finds what branch coverage missed.
\end{enumerate}
\end{workedbox}

\Needspace{4\baselineskip}
\Needspace{4\baselineskip}
\begin{notebox}\textbf{Trap:} Path coverage is evidence of execution, not proof of correctness.\end{notebox}

\Needspace{8\baselineskip}
\begin{notebox}\textbf{Near miss:} A branch test with \ensuremath{-}2 nearly checks abs, but \ensuremath{-}1 reveals the wrong condition.\end{notebox}

\begin{recallbox}{Recall \ensuremath{\cdot} §9.4}
\Needspace{4\baselineskip}
\begin{enumerate}
\item What does path-complete mean?
\item What does coverage not guarantee?
\end{enumerate}

\nopagebreak[4]{\small\itshape Answers: see “9.4 Glass-box testing” under Answers at the end.}
\end{recallbox}
\end{keepsection}

\begin{keepsection}
\Needspace{10\baselineskip}
\section{9.5 Debug systematically}

\Needspace{7\baselineskip}
\begin{plainbox}\textbf{In plain words:} find how the unexpected value arose; do not guess at a fix.\end{plainbox}

\begin{equationcard}{Debugging loop}
\[ \begin{gathered} \texttt{observe \ensuremath{\to} hypothesis \ensuremath{\to} repeatable test} \end{gathered} \]
{\small Change one cause at a time.}
\end{equationcard}

\Needspace{11\baselineskip}
Observe evidence, form a hypothesis, and repeat a small experiment. Use a simple failing input. Print function entry, parameters and results. With bisection, inspect midway through a long calculation, then focus on the half containing the first divergence.

\begin{workedbox}{Worked example \ensuremath{\cdot} find a wrong total}
\Needspace{8\baselineskip}
\begin{enumerate}
\item Save the smallest failing input.
\item Inspect an intermediate value halfway through.
\item If wrong, inspect the earlier half; otherwise inspect the later half.
\item Repeat to locate the first faulty operation.
\item Change one thing and rerun this case and regression tests.
\end{enumerate}
\end{workedbox}

\Needspace{4\baselineskip}
\Needspace{5\baselineskip}
\begin{notebox}\textbf{Trap:} Changing several lines before rerunning the case destroys useful evidence.\end{notebox}

\Needspace{9\baselineskip}
\begin{notebox}\textbf{Near miss:} A print at the final result shows failure but not where it began; inspect intermediate values.\end{notebox}

\begin{recallbox}{Recall \ensuremath{\cdot} §9.5}
\Needspace{4\baselineskip}
\begin{enumerate}
\item Why ask how the result arose?
\item What makes a debugging experiment useful?
\end{enumerate}

\nopagebreak[4]{\small\itshape Answers: see “9.5 Debug systematically” under Answers at the end.}
\end{recallbox}
\end{keepsection}

\begin{keepsection}
\Needspace{10\baselineskip}
\section{9.6 Exceptions and handlers}

\Needspace{7\baselineskip}
\begin{plainbox}\textbf{In plain words:} Python raises an exception when an operation cannot proceed.\end{plainbox}

\begin{equationcard}{Handler match}
\[ \begin{gathered} \texttt{raised type \ensuremath{\in} handled types} \end{gathered} \]
{\small Otherwise the exception propagates.}
\end{equationcard}

\Needspace{8\baselineskip}
\begin{itemize}
\item \texttt{IndexError}: index outside a sequence.
\item \texttt{NameError}: a name is not defined.
\item \texttt{TypeError}: an operation received an inappropriate type.
\item \texttt{ValueError}: a suitable type carries an invalid value.
\item \texttt{AttributeError}: an attribute lookup failed.
\end{itemize}

\Needspace{9\baselineskip}
A \texttt{try} body runs until an exception; a matching specific \texttt{except} handles it. Unmatched exceptions propagate.

\begin{workedbox}{Worked example \ensuremath{\cdot} safe conversion}
\texttt{try: int(text)} may raise \texttt{ValueError}; \texttt{except ValueError} can report a correction, and \texttt{else} runs only after success. Keep try narrow so unrelated errors remain visible.
\end{workedbox}

\Needspace{4\baselineskip}
\Needspace{5\baselineskip}
\begin{notebox}\textbf{Trap:} A bare except can hide a TypeError while handling a conversion error.\end{notebox}

\Needspace{9\baselineskip}
\begin{notebox}\textbf{Near miss:} Catching ValueError around the whole report almost works, but can mislabel a later ValueError; keep try narrow.\end{notebox}

\begin{recallbox}{Recall \ensuremath{\cdot} §9.6}
\Needspace{4\baselineskip}
\begin{enumerate}
\item How do TypeError and ValueError differ?
\item Why avoid bare except?
\end{enumerate}

\nopagebreak[4]{\small\itshape Answers: see “9.6 Exceptions and handlers” under Answers at the end.}
\end{recallbox}
\end{keepsection}

\begin{keepsection}
\Needspace{10\baselineskip}
\section{9.7 Raise or recover deliberately}

\Needspace{7\baselineskip}
\begin{plainbox}\textbf{In plain words:} catch an error only when there is a sensible next action.\end{plainbox}

\begin{equationcard}{Error policy}
\[ \begin{gathered} \texttt{recover only when the next action is defined} \end{gathered} \]
{\small Otherwise raise or re-raise a descriptive exception.}
\end{equationcard}

\Needspace{12\baselineskip}
Raise a descriptive exception if a function cannot meet its contract. Avoid magic error values that resemble valid results. A handler may notify, stop, apply a documented default, or re-raise. For empty grades, choose whether to reject or define a default average. Finally runs after success or failure, making it suitable for cleanup.

\begin{workedbox}{Worked example \ensuremath{\cdot} empty grades}
\Needspace{5\baselineskip}
\begin{enumerate}
\item sum(grades)/len(grades) raises ZeroDivisionError for an empty list.
\item One policy rejects the list; another defines an average of 0.0.
\item Document the choice so callers do not treat a default as measured data.
\end{enumerate}
\end{workedbox}

\Needspace{4\baselineskip}
\Needspace{4\baselineskip}
\begin{notebox}\textbf{Trap:} Returning a special number for failure can make it look like a valid result.\end{notebox}

\Needspace{8\baselineskip}
\begin{notebox}\textbf{Near miss:} Returning 0.0 for no grades is only valid if that policy is explicit.\end{notebox}

\begin{recallbox}{Recall \ensuremath{\cdot} §9.7}
\Needspace{4\baselineskip}
\begin{enumerate}
\item When re-raise?
\item What is finally for?
\end{enumerate}

\nopagebreak[4]{\small\itshape Answers: see “9.7 Raise or recover deliberately” under Answers at the end.}
\end{recallbox}
\end{keepsection}

\begin{keepsection}
\Needspace{10\baselineskip}
\section{9.8 Assertions and invariants}

\Needspace{7\baselineskip}
\begin{plainbox}\textbf{In plain words:} assert says an internal condition should always hold here.\end{plainbox}

\begin{equationcard}{Assertion}
\[ \begin{gathered} \texttt{assert condition} \end{gathered} \]
{\small Checks an internal assumption during development.}
\end{equationcard}

\Needspace{6\baselineskip}
\texttt{assert len(grades) \ensuremath{>} 0} raises AssertionError when false in a normal run. Assertions can check invariants and expected internal results during development. They supplement testing.

\Needspace{10\baselineskip}
Python may omit assertions under optimization (\texttt{-O}). Do not use them for required behavior or user-input validation; validate external data explicitly.

\begin{workedbox}{Worked example \ensuremath{\cdot} internal versus external}
\Needspace{5\baselineskip}
\begin{enumerate}
\item Assert an internal count remains non-negative after an algorithm step.
\item Explicitly reject invalid user input and display an error.
\item The first checks a programmer assumption; the second is required behavior.
\end{enumerate}
\end{workedbox}

\Needspace{4\baselineskip}
\Needspace{5\baselineskip}
\begin{notebox}\textbf{Trap:} An assertion is not a substitute for checking user data.\end{notebox}

\Needspace{9\baselineskip}
\begin{notebox}\textbf{Near miss:} assert age \ensuremath{>}= 18 looks like validation, but -O can remove it; use an explicit check.\end{notebox}

\begin{recallbox}{Recall \ensuremath{\cdot} §9.8}
\Needspace{4\baselineskip}
\begin{enumerate}
\item What is an invariant?
\item Why not assert external input?
\end{enumerate}

\nopagebreak[4]{\small\itshape Answers: see “9.8 Assertions and invariants” under Answers at the end.}
\end{recallbox}
\end{keepsection}

\begin{keepsection}
\Needspace{4\baselineskip}
\section{Decision map}
\end{keepsection}

\begin{keepsection}
\Needspace{11\baselineskip}
\section{Derive it cold}

\Needspace{8\baselineskip}
\begin{enumerate}
\item Design square-root test partitions. \hfill{\small\color{FAUblue}§9.3}
\item Trace the abs(-1) boundary defect. \hfill{\small\color{FAUblue}§9.4}
\item Write a specific try/except/else. \hfill{\small\color{FAUblue}§9.6}
\item Choose an empty-grade policy. \hfill{\small\color{FAUblue}§9.7}
\item Separate an invariant assertion from input validation. \hfill{\small\color{FAUblue}§9.8}
\end{enumerate}
\end{keepsection}

\begin{keepsection}
\Needspace{13\baselineskip}
\section{Explain it}

\Needspace{10\baselineskip}
\begin{enumerate}
\item Why write expected results before testing?
\item Why retain a regression test after fixing a defect?
\item Why doesn’t path coverage prove correctness?
\item Why use a small failing input?
\item Why keep try narrow?
\item Why avoid magic error values?
\end{enumerate}

\nopagebreak[4]{\small\itshape Answers: see “Explain it” under Answers at the end.}
\end{keepsection}

\clearpage

\begin{keepsection}
\Needspace{9\baselineskip}
\section{Mixed practice}

\Needspace{6\baselineskip}
20 problems. Try each before opening the matching worked answer.

\begin{enumerate}
\item State a contract for is\_bigger(x,y).
\item Which test level checks one function?
\item Name input classes for a first-item function on lists.
\item For n \ensuremath{<} 10, give boundary values on both sides.
\item Why can branch coverage miss the abs(-1) defect?
\item Describe bisection debugging.
\item Why preserve the smallest failing input?
\item What does int('3.5') raise?
\item What does [1,2][4] raise?
\item When raise instead of return -1?
\item Why avoid a bare except?
\item When does try/else execute else?
\item What belongs in finally?
\item What does a false assert do in a normal run?
\item Why not assert an access rule?
\item For eps\ensuremath{>}0, name an invalid boundary and a nearby valid value.
\item Why retain a regression test?
\item An empty grade list causes division by zero. Name a policy.
\item Should abs test -1 if 2 and -2 already ran?
\item What makes debugging systematic?
\end{enumerate}

\nopagebreak[4]{\small\itshape Answers: see “Mixed practice” under Answers at the end.}
\end{keepsection}

\begin{keepsection}
\Needspace{5\baselineskip}
\section{Exam-style problems}

Allow about 20 minutes for each problem.

\begin{keepsection}
\Needspace{13\baselineskip}
\subsection{Problem 1 \ensuremath{\cdot} Test a square-root function}

\Needspace{10\baselineskip}
Given x\ensuremath{\geq}0, eps\ensuremath{>}0, and the promise x\ensuremath{-}eps\ensuremath{\leq}r²\ensuremath{\leq}x+eps, design a test set, explain coverage limits, and describe debugging a failure.

\Needspace{7\baselineskip}
\begin{enumerate}
\item Partition the inputs.
\item State a boundary and extreme-value check.
\item Explain the path-coverage limitation.
\item Describe a debugging strategy.
\end{enumerate}

\nopagebreak[4]{\small\itshape Answers: see “Exam-style problems” under Answers at the end.}
\end{keepsection}

\begin{keepsection}
\Needspace{13\baselineskip}
\subsection{Problem 2 \ensuremath{\cdot} Grade report and errors}

\Needspace{10\baselineskip}
A class report computes each student's average. Choose behavior for empty grade lists, describe tests, and separate exception handling from assertions.

\Needspace{7\baselineskip}
\begin{enumerate}
\item State a policy for empty grades.
\item Name unit and regression tests.
\item Name an integration test.
\item Explain which checks should be assertions.
\end{enumerate}

\nopagebreak[4]{\small\itshape Answers: see “Exam-style problems” under Answers at the end.}
\end{keepsection}
\end{keepsection}

\begin{keepsection}
\Needspace{10\baselineskip}
\section{Slips in the notes}

\Needspace{7\baselineskip}
\begin{itemize}
\item Slides 21–22 use bare except before recommending specific handlers. This guide uses the specific form to keep unrelated failures visible.
\item Slide 33 shows assert as input validation. Python may remove assert with \texttt{-O}; explicitly validate user input.
\item Slide 24 lists silent failure but warns against it; an unexplained default can hide an error.
\end{itemize}
\end{keepsection}

\begin{keepsection}
\Needspace{7\baselineskip}
\section{Self-check}

\begin{enumerate}
\item $\mdlgwhtsquare$\  I state assumptions and promises (1).
\item $\mdlgwhtsquare$\  I distinguish unit, regression and integration (2, 17).
\item $\mdlgwhtsquare$\  I choose input classes and boundaries (3, 4, 16).
\item $\mdlgwhtsquare$\  I explain coverage limits (5, 19).
\item $\mdlgwhtsquare$\  I debug with repeatable evidence (6, 7, 20).
\item $\mdlgwhtsquare$\  I identify exception types (8, 9).
\item $\mdlgwhtsquare$\  I use specific handlers, else and finally (11, 12, 13).
\item $\mdlgwhtsquare$\  I choose an error policy (10, 18).
\item $\mdlgwhtsquare$\  I separate assertions from validation (14, 15).
\item $\mdlgwhtsquare$\  I distinguish exceptions from assertions (11, 14, 15).
\item $\mdlgwhtsquare$\  I derive the five cold problems.
\item $\mdlgwhtsquare$\  I solve both exam problems.
\end{enumerate}
\end{keepsection}


\clearpage
\section{Answers}
\begin{keepsection}\subsection{Start from memory}

\Needspace{9\baselineskip}
\begin{enumerate}
\item Testing compares actual behavior with a specification; debugging investigates a mismatch’s cause. \hfill{\small\color{FAUblue}§9.1}
\item Boundaries expose off-by-one and special-size errors. \hfill{\small\color{FAUblue}§9.3}
\item No; a path may run but compute the wrong result. \hfill{\small\color{FAUblue}§9.4}
\item ValueError: the type is suitable but the value is not an integer literal. \hfill{\small\color{FAUblue}§9.6}
\item No. Assertions may be removed under optimization. \hfill{\small\color{FAUblue}§9.8}
\end{enumerate}

\end{keepsection}
\begin{keepsection}\subsection{9.1 Defensive programming}

\Needspace{4\baselineskip}
\begin{enumerate}
\item Input assumptions and the promised result. \hfill{\small\color{FAUblue}§9.1}
\item To test and debug pieces independently. \hfill{\small\color{FAUblue}§9.1}
\end{enumerate}

\end{keepsection}
\begin{keepsection}\subsection{9.2 Test levels}

\Needspace{4\baselineskip}
\begin{enumerate}
\item A previously fixed defect. \hfill{\small\color{FAUblue}§9.2}
\item Whether components work together. \hfill{\small\color{FAUblue}§9.2}
\end{enumerate}

\end{keepsection}
\begin{keepsection}\subsection{9.3 Black-box testing}

\Needspace{4\baselineskip}
\begin{enumerate}
\item Boundary defects can be missed by ordinary examples. \hfill{\small\color{FAUblue}§9.3}
\item They can follow the specification without implementer bias. \hfill{\small\color{FAUblue}§9.3}
\end{enumerate}

\end{keepsection}
\begin{keepsection}\subsection{9.4 Glass-box testing}

\Needspace{4\baselineskip}
\begin{enumerate}
\item Every potential path was tested at least once. \hfill{\small\color{FAUblue}§9.4}
\item That every boundary or specification case is correct. \hfill{\small\color{FAUblue}§9.4}
\end{enumerate}

\end{keepsection}
\begin{keepsection}\subsection{9.5 Debug systematically}

\Needspace{4\baselineskip}
\begin{enumerate}
\item It focuses on the sequence of events and values. \hfill{\small\color{FAUblue}§9.5}
\item It is small, repeatable, and tests one hypothesis. \hfill{\small\color{FAUblue}§9.5}
\end{enumerate}

\end{keepsection}
\begin{keepsection}\subsection{9.6 Exceptions and handlers}

\Needspace{4\baselineskip}
\begin{enumerate}
\item Wrong operand type versus invalid value. \hfill{\small\color{FAUblue}§9.6}
\item It can hide unrelated programming errors. \hfill{\small\color{FAUblue}§9.6}
\end{enumerate}

\end{keepsection}
\begin{keepsection}\subsection{9.7 Raise or recover deliberately}

\Needspace{4\baselineskip}
\begin{enumerate}
\item When the function cannot safely recover. \hfill{\small\color{FAUblue}§9.7}
\item Cleanup regardless of success or failure. \hfill{\small\color{FAUblue}§9.7}
\end{enumerate}

\end{keepsection}
\begin{keepsection}\subsection{9.8 Assertions and invariants}

\Needspace{4\baselineskip}
\begin{enumerate}
\item A condition expected to remain true. \hfill{\small\color{FAUblue}§9.8}
\item Optimization can remove assertions. \hfill{\small\color{FAUblue}§9.8}
\end{enumerate}

\end{keepsection}
\begin{keepsection}\subsection{Explain it}

\Needspace{10\baselineskip}
\begin{enumerate}
\item Keeps expectation separate from observed behavior. \hfill{\small\color{FAUblue}§9.2}
\item Prevents a fixed defect returning. \hfill{\small\color{FAUblue}§9.2}
\item Coverage measures execution, not correctness. \hfill{\small\color{FAUblue}§9.4}
\item Removes irrelevant behavior. \hfill{\small\color{FAUblue}§9.5}
\item Prevents swallowing unrelated errors. \hfill{\small\color{FAUblue}§9.6}
\item A caller may confuse one with a valid result. \hfill{\small\color{FAUblue}§9.7}
\end{enumerate}

\end{keepsection}
\begin{keepsection}\subsection{Mixed practice}

\begin{enumerate}
\item Assume integer inputs. Return True exactly when x\ensuremath{>}y; otherwise return False. \hfill{\small\color{FAUblue}§9.1}
\item Unit testing checks an isolated function or component. \hfill{\small\color{FAUblue}§9.2}
\item Test a singleton, a longer list, and the empty list if its behavior is specified. \hfill{\small\color{FAUblue}§9.3}
\item Use 9 and 10; they lie immediately on opposite sides of the decision. \hfill{\small\color{FAUblue}§9.3}
\item The tests may take both branches at other values while missing the incorrect boundary. \hfill{\small\color{FAUblue}§9.4}
\item Inspect a midpoint, then focus on the half where the first value diverges. \hfill{\small\color{FAUblue}§9.5}
\item It makes the defect reproducible and removes unrelated behavior. \hfill{\small\color{FAUblue}§9.5}
\item ValueError: the text is not an integer literal. \hfill{\small\color{FAUblue}§9.6}
\item IndexError: index 4 is outside the list. \hfill{\small\color{FAUblue}§9.6}
\item When no valid result can be produced or -1 may be mistaken for a valid result. \hfill{\small\color{FAUblue}§9.7}
\item It can hide unrelated errors; catch the expected exception specifically. \hfill{\small\color{FAUblue}§9.6}
\item Only when the try body completes without raising an exception. \hfill{\small\color{FAUblue}§9.6}
\item Cleanup that must run whether the operation succeeds or fails. \hfill{\small\color{FAUblue}§9.7}
\item It raises AssertionError. \hfill{\small\color{FAUblue}§9.8}
\item Python can omit assertions under optimization, so required checks must be explicit. \hfill{\small\color{FAUblue}§9.8}
\item eps=0 is excluded; a small positive epsilon is valid. \hfill{\small\color{FAUblue}§9.3}
\item It catches the fixed defect if a later change reintroduces it. \hfill{\small\color{FAUblue}§9.2}
\item Reject with ValueError, or document an intentional default such as 0.0. \hfill{\small\color{FAUblue}§9.7}
\item Yes. -1 is the boundary that exposes the faulty condition. \hfill{\small\color{FAUblue}§9.4}
\item Use evidence, form one hypothesis, and repeat a controlled test. \hfill{\small\color{FAUblue}§9.5}
\end{enumerate}

\end{keepsection}
\begin{keepsection}\subsection{Exam-style problems}

\Needspace{10\baselineskip}
\begin{enumerate}
\item Test zero, perfect squares, sub-unit values, irrational roots, and very small and large values. \hfill{\small\color{FAUblue}§9.3}
\item Check zero and extremes; evaluate the promised inequality instead of requiring one exact approximation. \hfill{\small\color{FAUblue}§9.3}
\item Exercise branches and loop exits, but keep boundary tests because execution coverage does not prove correctness. \hfill{\small\color{FAUblue}§9.4}
\item Keep the failing input and bisect intermediate values to find the first divergence. \hfill{\small\color{FAUblue}§9.5}
\end{enumerate}

\end{keepsection}
\begin{keepsection}\subsection{Exam-style problems}

\Needspace{8\baselineskip}
\begin{enumerate}
\item Choose and document rejection or a defined default. \hfill{\small\color{FAUblue}§9.7}
\item Test ordinary and empty inputs; preserve the divide-by-zero case as a regression test. \hfill{\small\color{FAUblue}§9.2, §9.7}
\item Test a full class report with ordinary and empty records. \hfill{\small\color{FAUblue}§9.2}
\item Validate outside input explicitly; use assertions for internal invariants. \hfill{\small\color{FAUblue}§9.8}
\end{enumerate}

\end{keepsection}

\vfill{\small Based on the supplied “09\_testing\_debugging.pdf”, rewritten in my own words. See the Python exception tutorial and assert reference. Not an official FAU publication. Made by \textbf{Thanate Amornjiraporn (Focus)}.}

\end{document}
